\chapter{球面、双曲空间与\(k\)的含义}

让我们做一点数学，一些简单的数学。任意维数中的球面都是特殊的几何：它们是均匀的，而这里的均匀意味着常曲率。正曲率首先出现，然后是均匀负曲率，再之后是与爱因斯坦方程的联系。

球面宇宙意味着\emph{空间}是球形的，而不是时空。还有另一种可能性，即均匀负曲率。把这两种情形并置在一起，才能让广义相对论的作用变得清楚。

\section{球面与均匀几何}

一维球面就是一个圆。描述它的最简单方式是把它嵌入一个二维辅助空间，并写出
\begin{equation}
x^2+y^2=R^2.
\end{equation}
单位一维球面是
\begin{equation}
x^2+y^2=1.
\end{equation}

二维球面以同样的方式描述，现在使用三个辅助坐标：
\begin{equation}
x^2+y^2+z^2=R^2,
\end{equation}
单位情形为
\begin{equation}
x^2+y^2+z^2=1.
\end{equation}

这些坐标只是一根拐杖。我们关心的几何并不是周围的\(x,y,z\)空间，而是曲面本身。生活在一维球面上的生物生活在一个圆上；它无法进入内部或外部。同样的想法适用于二维球面。球面的填充部分和球面的外部只属于辅助图像。

同样的构造给出三维球面。引入多一个坐标，比如\(w\)，并写出
\begin{equation}
x^2+y^2+z^2+w^2=R^2,
\end{equation}
单位情形为
\begin{equation}
x^2+y^2+z^2+w^2=1.
\end{equation}
三维球面更难可视化，因为此时拐杖将是四维的，但数学并没有难多少。

标准记号是用\(\Omega_n\)表示\(n\)维球面。因此\(\Omega_1\)是一维球面，\(\Omega_2\)是二维球面，\(\Omega_3\)是三维球面。

\section{球面上的度规}

更有用的描述是内禀的。人们用生活在球面本身的坐标来描述球面。

对于单位一维球面，一个坐标\(\phi\)就足够了，其中\(0\leq \phi \leq 2\pi\)。其度规为
\begin{equation}
ds^2=d\phi^2.
\end{equation}
这也写作
\begin{equation}
d\Omega_1^2=d\phi^2.
\end{equation}

单位二维球面被描述为一串一维球面。一个角\(\phi\)仍然绕着每个圆转，但第二个角\(\theta\)从球面的一端延伸到另一端。这第二个角从\(0\)到\(\pi\)。沿\(\theta\)方向的运动贡献\(d\theta^2\)。绕\(\phi\)方向的运动发生在一个半径为\(\sin\theta\)的圆上。因此
\begin{equation}
d\Omega_2^2=d\theta^2+\sin^2\theta\,d\phi^2
\end{equation}
或
\begin{equation}
d\Omega_2^2=d\theta^2+\sin^2\theta\,d\Omega_1^2.
\end{equation}

三维球面以同样的方式得到，现在是一串二维球面。引入另一个坐标\(\alpha\)。于是
\begin{equation}
d\Omega_3^2=d\alpha^2+\sin^2\alpha\,d\Omega_2^2.
\end{equation}
展开后，这就是
\begin{equation}
d\Omega_3^2=d\alpha^2+\sin^2\alpha\left(d\theta^2+\sin^2\theta\,d\phi^2\right).
\end{equation}
这里\(\alpha\)从一个零尺寸端延伸到另一端，所以\(0\leq \alpha \leq \pi\)。\editorialnote{在00：17：24处，音频短暂听起来像``零到二{\proseunicodefont π}，''但周围的构造以及显式的\(\sin^2\alpha\)因子要求新坐标从低维球面族的零尺寸一端延伸到另一端。}

一般的递推形式是
\begin{equation}
d\Omega_n^2=d\alpha^2+\sin^2\alpha\,d\Omega_{n-1}^2.
\end{equation}

一旦采用这种观点，嵌入图像就可以被遗忘。一维球面由一个坐标描述，二维球面由两个坐标描述，三维球面由三个坐标描述，依此类推。几何自身独立存在。

这种几何之所以特殊，是因为它是齐性的。人们可以通过把\(\sin^2\theta\)替换为某个仍在两端为零的其他函数来变形度规。这样就会得到某个变形的闭合曲面，例如椭球面。但它将不再是均匀的。球面是齐性的，因为每一点都与任何其他点相同。

如果想要半径为\(A\)的球面而不是单位半径，就把单位度规乘以\(A^2\)：
\begin{equation}
ds^2=A^2\,d\Omega_3^2.
\end{equation}
平方是必要的，因为度规测量的是距离的平方。

一个简短的题外话涉及零维球面。向后外推一步给出
\begin{equation}
x^2=1,
\end{equation}
其解是两个点\(x=\pm 1\)。所以单位零维球面就是两个点。

\section{球形的FRW宇宙}

现在关心的不只是空间，而是时空。相关的图像是一个半径随时间变化的三维球面。

有一个警告始终重要：时间是不同的。人们不应把膨胀宇宙想成一个高维球面，而让时间扮演又一个嵌入坐标的角色。

考虑到这一点，时空间隔为
\begin{equation}
d\tau^2=dt^2-a(t)^2\,d\Omega_3^2.
\end{equation}
这就是空间上球形的宇宙的FRW度规。与静态球面相比，唯一的新要素是半径变成了一个依赖于时间的函数\(a(t)\)。

选择坐标使我们位于一个极点是方便的。那个点并没有什么特殊之处；这只是一个约定。于是\(\alpha\)可以被看作单位球面上离开我们的径向坐标。它不是以千米为单位的径向距离。它是单位球面上的无量纲角量度，但它在几何中扮演径向距离的角色。

除了星系和恒星这样的局部团块之外，这就是空间球对称FRW宇宙的全部几何。它的空间曲率为正，符号约定写作
\begin{equation}
k=+1.
\end{equation}
这本身并不是``加一''。它仅仅意味着正空间曲率。

\section{从球面向外看}

当一个人从球面上的某一点向外看时，这样的球面看起来是什么样子？这正是球极投影登场之处。

取一个与无限平面在南极点相切的二维球面。在球面上选一个点。从北极点经过该点画一条直线，并延长它直到它与平面相交。这样，球面上的每个点都被映射到平面上的一个点。南极点落在平面的原点。北极点是例外点：当一个人从不同方向趋近它时，它的像以依赖于方向的方式奔向无穷远。因此人们把它当作无穷远点来处理。

如果把相等的``大陆''放在球面上不同的位置，它们在平面上的像会随着人们向北极点移动而变得更大。因此，位于南极点的观察者会看到遥远天体所张的角度比它们在平坦空间中张的角度更大。在正曲率中，一个已知天体相对于欧几里得几何看起来太大了。

这就是基本的观测要点。曲率是从距离、大小和角大小之间的关系推断出来的。星系本身对于非常深的距离而言并不是好的标准尺，但逻辑恰恰就是如此。

\begin{classroomqa}
\lecturetimestamp{00:45:54}
\audiencequestion{你是否也必须知道距离，而不仅仅是角大小？}
\lecturerresponse{是的。仅有角度是不够的。你需要角度和距离。}
\end{classroomqa}

\begin{classroomqa}
\lecturetimestamp{00:46:18}
\audiencequestion{球面上的圆会投影为圆吗？}
\lecturerresponse{是的。圆投影为圆。这是球极投影的特殊特征之一。}
\end{classroomqa}

\begin{classroomqa}
\lecturetimestamp{00:47:02}
\audiencequestion{一个大陆在投影下会保持相同的形状吗？}
\lecturerresponse{这是两个不同的陈述。一个是圆投影为圆。另一个是较弱的陈述：足够小的形状保持其局部形状，因为角度被保持。该映射是共形的：它保持角度，尽管不保持相对大小。}
\end{classroomqa}

\section{均匀负曲率}

均匀负曲率的情形是通过其度规引入的。

球面的负曲率类比是双曲空间。按照与之前相同的递归精神，其度规通过把\(\sin\alpha\)替换为\(\sinh\alpha\)而得到：
\begin{equation}
d\mathcal H_n^2=d\alpha^2+\sinh^2\alpha\,d\Omega_{n-1}^2.
\end{equation}
低维角部分仍然是单位球面度规\(d\Omega_{n-1}^2\)。

\begin{classroomqa}
\lecturetimestamp{00:55:51}
\audiencequestion{在度规中，那个低维微分仍然是球面的吗？}
\lecturerresponse{是的。那个低维部分就是普通的单位球面度规。变化在于它前面的\(\sinh^2\alpha\)因子。}
\end{classroomqa}

双曲正弦是
\begin{equation}
\sinh\alpha=\frac{e^\alpha-e^{-\alpha}}{2}.
\end{equation}
负号很重要；在讲座中它立刻被一名学生纠正了。

坐标的范围也有一个关键区别。对于球面，新坐标从一端到另一端，并在\(\pi\)处停止。这里，
\begin{equation}
0\leq \alpha<\infty.
\end{equation}
普通正弦上升然后回到零。双曲正弦在\(\alpha=0\)附近以相同方式开始，但随后它指数增长。

这意味着我们周围的低维球面随着我们向外移动而变得越来越大。在这种几何中，半径增大处的天空不会先变大然后缩小。它持续变大。

在三个空间维度中，相应的FRW形式是
\begin{equation}
d\tau^2=dt^2-a(t)^2\,d\mathcal H_3^2,
\end{equation}
而符号标记是
\begin{equation}
k=-1.
\end{equation}

\section{圆盘图景及其含义}

负曲率在由埃舍尔使之闻名的圆盘图景中变得直观，尤其是\emph{极限圆第四号}。与球极投影的类比是有用的，但行为是相反的。

对于球面，投影看起来是无限的，而真实几何是有限的。对于负曲率情形，圆盘图像看起来是有限的，而真实几何是无限的。相等的``大陆''或``星系''在向边界移动时显得越来越小，但几何上它们的大小全都相同。如果通过数穿过多少等份来测量距离，那么边界就在无限远处。

因此，负曲率中的一个已知物体相对于平坦空间看起来太小。它张开的角太小。

几何仍然是均匀的。尽管图像使外侧的部分看起来极小，但任何被移到另一位置的观测者都会看到相同的内禀几何。

\begin{classroomqa}
\lecturetimestamp{01:06:42}
\audiencequestion{处于这些圆中任意一个里的观测者会看到类似的图样吗？}
\lecturerresponse{是的，完全相同的东西。空间是均匀的。}
\end{classroomqa}

\begin{classroomqa}
\lecturetimestamp{01:07:03}
\audiencequestion{在球面情形中，远处的东西似乎太大，而在这里它们似乎太小？}
\lecturerresponse{是的。在正曲率中，一个已知物体相对于平坦空间张开的角太大。在负曲率中，它张开的角太小。}
\end{classroomqa}

\begin{classroomqa}
\lecturetimestamp{01:08:09}
\audiencequestion{这种球极投影式投影的意义在于天文学家能够从观测中重建它吗？}
\lecturerresponse{是的。天文学家实际上就是一位测量员，测量天空中的三角形。就眼下这个几何要点而言，可以暂且搁置时间依赖性，专注于角度和距离。}
\end{classroomqa}

还有一个有用的局部图像。负曲率的一小块区域可以画成普通三维空间中的一个鞍面。但只有小块区域能这样画；整个双曲平面不能作为普通三维欧几里得空间中的一个良好普通曲面整体嵌入。

\section{平坦情形}

在正曲率情形与负曲率情形之间，是平坦空间。可以把它看作一种极限，但只需把它想成普通的平坦欧几里得空间就足够了。

它的 FRW 形式是
\begin{equation}
d\tau^2=dt^2-a(t)^2\left(dx^2+dy^2+dz^2\right),
\end{equation}
其中
\begin{equation}
k=0.
\end{equation}

这就是三种均匀 FRW 几何：
\begin{align}
k=+1 &: \text{positive spatial curvature},\\
k=0 &: \text{flat space},\\
k=-1 &: \text{negative spatial curvature}.
\end{align}

\begin{classroomqa}
\lecturetimestamp{01:15:07}
\audiencequestion{如果把几何投影到一个平面上并追踪一个信号，投影出的光速会随位置变化吗？}
\lecturerresponse{是的，在投影上会。在球面投影上，越往外它看起来会越来越快；在负曲率圆盘上，靠近边缘时它看起来会越来越慢。这种变化属于图像，而不属于底层物理。}
\end{classroomqa}

\begin{classroomqa}
\lecturetimestamp{01:17:05}
\audiencequestion{人们实际上曾用这种方式观察远处物体来确定平坦性吗？}
\lecturerresponse{是的，但真正使用的是角度信息。人们并不是通过试图测量光在投影上看起来移动得多快来确定平坦性。}
\end{classroomqa}

\section{回到宇宙学方程}

休息之后，是已经研究过的宇宙学方程与爱因斯坦方程之间的联系。此时黑板上仍然保留着平坦空间度规的片段。

\begin{figure}[htbp]
\centering
\includegraphics[width=0.72\textwidth]{lecture_05_figure_02.png}
\caption{平坦 FRW 度规片段。}
\lectureframe{01:18:18}
\end{figure}

首先是一个非常简短的牛顿力学回顾。把我们放在中心。观察一个按照哈勃定律远离的星系。根据牛顿关于球对称分布的定理，只有半径与该星系距离相匹配的球面内部的物质对其所受的力有贡献。于是有三种可能：该星系以高于逃逸速度运动、低于逃逸速度运动，或恰好以逃逸速度运动。

接下来回顾的是，对于普通非相对论物质，密度按质量除以\(a^3\)变化；对于辐射，它按\(1/a^4\)变化；对于真空能，它不随球面的变化而变化。

在先前的牛顿处理中，宇宙学方程有一个密度项，并且如果人们不恰好处于逃逸速度情形，还有一个正比于\(1/a^2\)的额外常数项。先前的讨论大多把该项设为零。从经验上看它很小，但从概念上说没有理由把它省略。

\section{爱因斯坦方程与\(00\)分量}

下一步是爱因斯坦场方程。左侧是几何，由时空曲率构建。右侧是能量-动量张量，即引力场的源。在讲座的归一化中，人们写作
\begin{equation}
R^{\mu\nu}-\frac{1}{2}g^{\mu\nu}R=4\pi G\,T^{\mu\nu}.
\end{equation}
\editorialnote{在可见的板书行中，因子\(\tfrac12\)无法辨认，但在01:33:36附近，口头更正将其恢复。此处的方程保留了该更正，同时保持了讲座的\(4\pi G\)归一化。}

这里并不需要全部十个独立分量。到这一步，已经假设了太多---均匀性、FRW形式以及单个依赖于时间的尺度因子---，以至于只剩下一个任意函数，即\(a(t)\)。自然选择的组分是时间-时间分量，因为\(T^{00}\)是能量密度\(\rho\)。

\begin{figure}[htbp]
\centering
\includegraphics[width=0.62\textwidth]{lecture_05_figure_05.png}
\caption{爱因斯坦场方程及其\(00\)分量。}
\lectureframe{01:30:54}
\end{figure}

因此焦点是
\begin{equation}
R^{00}-\frac{1}{2}g^{00}R=4\pi G\,\rho.
\end{equation}

右侧是能量密度。左侧是几何。

\begin{classroomqa}
\lecturetimestamp{01:39:09}
\audiencequestion{那么能量-动量张量的其他分量呢？}
\lecturerresponse{它们一般不为零。空间-空间分量可以表示压强。在物质主导的情况下，压强为零；在辐射主导的情况下，压强为能量密度的三分之一；对于真空能，人们写作\(P=w\rho\)，其中\(w=-1\)。}
\end{classroomqa}

\begin{classroomqa}
\lecturetimestamp{01:40:24}
\audiencequestion{那么，如果只有物质，\(T^{\mu\nu}\)的其他分量会消失吗？}
\lecturerresponse{对于无压物质，是的，压强分量消失。}
\end{classroomqa}

\section{爱因斯坦计算给出了什么}

完整的曲率计算并未在板上进行。重要的是将FRW度规代入爱因斯坦张量后所出现的结果。

一类项仅仅来自尺度因子依赖于时间这一事实。即使空间是平坦的，依赖于时间的尺度因子仍会给出时空曲率。该贡献正比于
\begin{equation}
\left(\frac{\dot a}{a}\right)^2.
\end{equation}

第二类项来自空间曲率本身。即使尺度因子在时间上冻结，弯曲空间仍会对时空曲率有贡献。对于简单的均匀弯曲几何，曲率与曲率半径的平方成反比，因此它贡献一个正比于
\begin{equation}
\frac{k}{a^2},
\end{equation}
的项，球面为正号，双曲情形为负号，平坦空间为零。

\begin{classroomqa}
\lecturetimestamp{01:36:53}
\audiencequestion{是否还应该有一个诸如曲率半径的因子？}
\lecturerresponse{那已经是\(1/a^2\)因子了。对于这些简单的均匀弯曲几何，曲率就是曲率半径平方的倒数。}
\end{classroomqa}

系数的记账是即兴完成的，并部分与先前的牛顿结果进行了核对。最终结果是\footnote{这个展示的弗里德曼形式整合了口头的、自我修正的系数讨论，以及在01:32:35--01:36:53附近与先前牛顿结果的比较。}
\begin{equation}
\left(\frac{\dot a}{a}\right)^2+\frac{k}{a^2}=\frac{8\pi G}{3}\rho,
\end{equation}
或等价地
\begin{equation}
\left(\frac{\dot a}{a}\right)^2=\frac{8\pi G}{3}\rho-\frac{k}{a^2}.
\end{equation}

这就是将三种FRW几何与先前的宇宙动力学联系起来的那个形式。

\begin{classroomqa}
\lecturetimestamp{01:37:52}
\audiencequestion{如果重新定义\(a\)，\(k\)可以取任意值吗？}
\lecturerresponse{在平坦情形下，\(a\)没有唯一的自然尺度，因此可以重新标度它，但\(k\)仍为零。在球面和负曲率情形下，存在一个自然的曲率尺度。如果取\(a(t)\)为该自然曲率半径，那么\(k\)就只是\(+1\)或\(-1\)。}
\end{classroomqa}

\begin{classroomqa}
\lecturetimestamp{01:41:00}
\audiencequestion{在最终方程中，\(k\)等于\(+1\)、\(-1\)还是\(0\)？}
\lecturerresponse{是的。对于球面它是\(+1\)，对于双曲情形它是\(-1\)，对于平坦空间它是\(0\)。}
\end{classroomqa}

牛顿的处理并没有把我们引入歧途；它给出了相同的基本方程。广义相对论添加的是解释。从牛顿的观点来看，额外的常数区分了高于逃逸速度、低于逃逸速度以及临界情形。从爱因斯坦的观点来看，相同的符号告诉我们空间本身的几何。

\begin{classroomqa}
\lecturetimestamp{01:42:46}
\audiencequestion{在引入宇宙学常数之前，这一切是如何成立的？}
\lecturerresponse{在宇宙学常数出现之前，逃逸速度行为与空间几何之间存在紧密关系。若\(k=+1\)，空间为正曲率，运动低于逃逸速度，因此宇宙可以先膨胀再收缩。若\(k=-1\)，空间为负曲率，运动高于逃逸速度，因此宇宙持续膨胀。若\(k=0\)，则处于临界的刀锋情形。}
\end{classroomqa}

因此，同一个符号\(k\)可以有两种解读：作为运动的动力学分类，以及作为空间曲率的符号。
